\subsection{\texorpdfstring{\(\mathcal{L}_{\mathfrak{S}}(\mathbf{E};\mathbf{F})\) 为 Hausdorff 的条件}{\textbackslash mathcal\{L\}\_\{\textbackslash mathfrak\{S\}\}(\textbackslash mathbf\{E\};\textbackslash mathbf\{F\}) 为 Hausdorff 的条件}}

命题 3. -- 设 E 和 F 是两个局部凸空间，其中 F 假设为 Hausdorff 的，并设 \(\mathfrak{S}\) 是 E 的有界子集的一个族。若 \(\mathfrak{S}\) 中诸集的并 A 在 E 中是完全的，则空间 \(\mathcal{L}_{\mathfrak{S}}(\mathrm{E};\mathrm{F})\) 是 Hausdorff 的。

设 \(u_0\) 是 \(\mathcal{L}(\mathrm{E};\mathrm{F})\) 中一个非零元素；由于 \(u_0\) 连续且 \(\mathrm{A}\) 在 \(\mathrm{E}\) 中是完全的，存在一点 \(x_0\) 在 \(\mathrm{A}\) 中，使得 \(u_0(x_0) \neq 0\)。由于 \(\mathrm{F}\) 是 Hausdorff 的，存在 0 的一个邻域 \(\mathrm{V}\)（在 \(\mathrm{F}\) 中）使得 \(u_0(x_0) \notin \mathrm{V}\)。设 \(\mathbf{M} \in \mathfrak{S}\) 满足 \(x_0 \in \mathbf{M}\)。于是集 \(\mathrm{U}\)（由所有 \(u \in \mathcal{L}(\mathrm{E};\mathrm{F})\) 中使 \(u(\mathbf{M}) \subset \mathbf{V}\) 者组成）是 \(\mathcal{L}(\mathrm{E};\mathrm{F})\) 中 0 的一个邻域，且我们有 \(u_0 \notin \mathrm{U}\)，从而 \(\mathcal{L}(\mathrm{E};\mathrm{F})\) 是 Hausdorff 的。

特别地，只要 F 是 Hausdorff 的，\(\mathcal{L}(\mathrm{E};\mathrm{F})\) 上的下列拓扑就都是 Hausdorff 的：简单收敛、紧收敛、预紧或紧凸收敛，以及有界收敛拓扑。

